\documentclass[10pt]{article}
\usepackage[margin=0.55in]{geometry}
\usepackage[T1]{fontenc}
\usepackage{lmodern}
\usepackage{enumitem}
\usepackage{booktabs}
\usepackage[hidelinks]{hyperref}
\usepackage{titlesec}\titlespacing*{\section}{0pt}{6pt}{3pt}
\setlist{nosep,leftmargin=1.2em}
\setlength{\parindent}{0pt}
\setlength{\parskip}{4pt}
\pagestyle{empty}

\begin{document}

{\Large\bfseries HAALSA Hackathon --- results brief}\hfill Sat Oct 24, 10 AM--5 PM ET

Zoom: \url{https://bostonu.zoom.us/j/96853036635} \quad ID 968 5303 6635 \quad passcode 421911.\\
Folder \texttt{research/haalsi-hackathon/}: untouched starter \texttt{HAALSI\_hackathon\_student.Rmd};\\
filled copy with all answers \texttt{HAALSI\_hackathon\_joshua.Rmd} / \texttt{.html}.

\section*{Pitch: Question 1 --- which intervention could cut cognitive impairment most?}

\textbf{Data.} HAALSI, rural Agincourt, South Africa: 5{,}059 adults aged 40+ at Wave 1, three
waves over about 6 years (12{,}092 person-waves with cognition). \textbf{Outcome}:
\texttt{objective\_cognition\_z} (orientation + immediate + delayed word recall, Wave 1 SD units)
and \texttt{low\_cognition} (below the Wave 1 20th percentile). \textbf{Confounders}: age, sex,
education (all come before the exposures). Wealth is added for hypertension and depression. BMI,
hypertension and depression are \emph{not} adjusted for in the wealth model because they may be
mechanisms.
\textbf{Methods}: linear regression (level), logistic regression (low cognition), linear mixed
model with random intercept and slope (change), then predicted prevalence under each intervention.

\begin{center}\small
\begin{tabular}{@{}lccc@{}}
\toprule
Exposure & Level (SD) & OR low cognition & Modelled change in low cognition\\
\midrule
Wealth, per 1 SD assets & $+0.11$ (0.08, 0.14) & 0.70 (0.63, 0.78) & $+0.5$ SD assets: 19.9\% $\to$ 17.8\% ($-2.2$ pts)\\
Hypertension, untreated & $+0.01$ ($-0.04$, 0.07) & 1.01 (0.84, 1.23) & \\
Hypertension, treated & $+0.16$ (0.10, 0.22) & 0.61 (0.50, 0.76) & treat all untreated: $-2.1$ pts\\
Depression, per 1 SD CES-D & $-0.11$ ($-0.13$, $-0.08$) & 1.26 (1.17, 1.36) & $-0.5$ SD if symptomatic: $-0.6$ pts\\
\bottomrule
\end{tabular}
\end{center}

\textbf{How to read it.}
\begin{itemize}
\item \textbf{Wealth}: graded --- richest vs poorest quintile $+0.31$ SD. A 0.5 SD rise in assets
  removes about 11\% of low-cognition cases. It does \emph{not} change the 6-year slope
  (0.002 SD/yr).
\item \textbf{Hypertension}: untreated looks the same as none. The whole ``benefit'' is treated
  vs untreated people, who differ in clinic access, health literacy and resources
  (healthy-user bias). This is the estimate we trust least.
\item \textbf{Depression}: linked to self-reported difficulty almost twice as strongly
  (0.17--0.20 SD) as to test scores ($-0.11$). Self-reports over-flag depressed people.
\end{itemize}
\textbf{Verdict}: financial security (cash transfer / Old Age Pension) is the most promising. It has
the largest modelled effect, the most credible mechanism and existing natural-experiment evidence
(Case 2001; Duflo 2003; Ardington et al.\ 2009), and it benefits the whole household.

\textbf{Limitations} (pick 1--2): (1) confounding by lifelong factors --- childhood
circumstances and cognitive ability shape both wealth and scores. (2) Healthy-user bias in the
treated-hypertension contrast. (3) Practice effects and dropout or death (about 1 in 5 by Wave 3):
scores \emph{rise} 0.08 SD/yr on retest, so ``decline'' analyses compare gains. (4) A short
cognitive test with ceiling effects.

\textbf{Implication / better study}: a regression discontinuity at pension-eligibility age 60, or a
randomized cash transfer with cognitive follow-up.

\section*{Backup findings (if the team picks another question)}
\begin{itemize}
\item \textbf{Q2 self-rated memory}: Spearman $\rho=-0.38$. ``Fair/poor'' catches only 37\% of
  low-cognition cases (specificity 0.86, PPV 0.40). Related, but a poor screen.
\item \textbf{Q3 diabetes}: objective cognition no different ($-0.07$, CI crosses 0), but worse
  self-rated memory ($+0.13$ SD). Only 249 treated --- imprecise.
\item \textbf{Q4 BMI}: underweight $-0.20$ SD, obese $+0.20$ SD, non-linear (spline $p<0.001$).
  Likely reverse causation (illness or poverty causes thinness). No moderation by hypertension
  ($p=0.45$).
\item \textbf{Q5 ADLs}: hypertension \emph{and} diabetes together $\approx$ doubles odds of any ADL
  difficulty (OR 2.02, 1.42--2.85). Either alone is not clearly raised. Interaction underpowered.
\item \textbf{Q6/Q7 change}: education and older age predict \emph{smaller practice gains} (a
  ceiling and regression-to-the-mean artefact, not harm). In the residual approach, wealth
  predicts outperforming one's expected score (OR 1.25/SD), but not in the slope model --- the two
  approaches disagree.
\end{itemize}

\textbf{Slide outline (5--6 slides)}: question \& why it matters $\to$ data + variables + why
those confounders $\to$ methods $\to$ the table above $\to$ verdict + limitations $\to$ better
study. Required citation on the final slide: \emph{HAALSI is sponsored by the National Institute on
Aging (5P01AG041710-05) and conducted by the Harvard Center for Population and Development Studies
in partnership with Wits University.}

\end{document}
